\rtmtight
\begin{tblr}{width=\textwidth,colspec={|Q[c,m,2.1cm]|X[l,m]|},hlines,vlines,hspan=minimal,rowsep=1.5pt,colsep=3pt,row{1}={bg=headgray,font=\bfseries}}
\SetCell{c,m}\hfil\logo\hfil & \textbf{POLITEKNIK NEGERI MALANG}\par\textbf{JURUSAN TEKNOLOGI INFORMASI}\par\textbf{PROGRAM STUDI: D4 TEKNIK INFORMATIKA} \\
\SetCell[c=2]{l} \textbf{RENCANA TUGAS MAHASISWA (RTM) DAN RUBRIK PENILAIAN} & \\
Nama Mata Kuliah & Praktikum Jaringan Komputer \\
Kode Mata Kuliah & RTI254006 \quad \textbf{sks / Jam perminggu:} 2 SKS/4 jam \quad \textbf{Semester:} 4 \\
Dosen Pengampu & Tim Pengajar Program Studi D4 TI \\
\end{tblr}
\assessmentblock{Tugas Jobsheet dan Laporan Praktikum}{Praktikum dan Laporan}{SCPMK0504-03301, SCPMK0504-03302, SCPMK0901-03303}{Mahasiswa menyelesaikan jobsheet praktikum dan menyusun laporan hasil konfigurasi/analisis sesuai Sub-CPMK dan konteks mata kuliah.}{Melakukan konfigurasi perangkat jaringan, menguji konektivitas, mendokumentasikan langkah-langkah, menganalisis hasil, dan menyusun laporan terstruktur.}{Jobsheet terisi, laporan praktikum, screenshot/dokumentasi konfigurasi, dan hasil pengujian.}{Ketepatan konfigurasi dan prosedur & 9 & Konfigurasi dilakukan sesuai prosedur, parameter jaringan tepat, dan konektivitas berhasil diverifikasi. & Konfigurasi sebagian besar benar, tetapi terdapat parameter yang kurang tepat atau verifikasi belum lengkap. & Konfigurasi salah, tidak mengikuti prosedur, atau konektivitas tidak dapat diverifikasi. \\ \\
Kualitas analisis dan troubleshooting & 7 & Masalah jaringan diidentifikasi dengan tepat, analisis sistematis, dan solusi terdokumentasi dengan jelas. & Analisis dilakukan tetapi kurang sistematis atau dokumentasi solusi belum lengkap. & Tidak ada analisis bermakna atau tidak mampu mengidentifikasi masalah jaringan. \\ \\
Kelengkapan laporan dan dokumentasi & 4 & Laporan memuat semua komponen (tujuan, langkah, hasil, kesimpulan) dan screenshot/dokumentasi lengkap. & Laporan tersedia tetapi beberapa komponen atau dokumentasi masih kurang lengkap. & Laporan tidak terstruktur atau dokumentasi sangat minim. \\}

\assessmentblock{Kuis Konsep dan Troubleshooting Jaringan}{Kuis}{SCPMK0504-03302, SCPMK0901-03303}{Mahasiswa menyelesaikan kuis konsep dan troubleshooting jaringan sesuai Sub-CPMK.}{Menjawab soal tertulis tentang konsep protokol, subnetting, routing, dan troubleshooting jaringan.}{Lembar jawaban kuis.}{Penguasaan konsep protokol jaringan & 4 & Jawaban menunjukkan pemahaman mendalam tentang protokol TCP/IP, ARP, routing, dan cara kerjanya. & Sebagian besar konsep dijelaskan benar, tetapi terdapat pemahaman yang kurang tepat. & Jawaban keliru atau hanya menghafal istilah tanpa memahami cara kerja protokol. \\ \\
Ketepatan perhitungan subnetting & 3.5 & Perhitungan subnet mask, network address, broadcast, dan host range dilakukan benar dan efisien. & Perhitungan subnetting sebagian besar benar, tetapi terdapat kesalahan minor pada satu atau dua langkah. & Perhitungan subnetting salah atau tidak dapat menyelesaikan soal subnetting. \\ \\
Analisis skenario troubleshooting & 2.5 & Skenario masalah jaringan diidentifikasi dengan tepat, root cause dianalisis, dan solusi yang diusulkan relevan. & Identifikasi masalah cukup tepat, tetapi analisis root cause atau solusi belum lengkap. & Tidak dapat mengidentifikasi masalah atau solusi yang diusulkan tidak relevan. \\}

\assessmentblock{UTS Praktikum Tengah Semester}{Ujian Praktikum}{SCPMK0504-03301, SCPMK0504-03302}{Mahasiswa menyelesaikan ujian praktikum mid-semester yang menguji kemampuan konfigurasi perangkat jaringan dan subnetting.}{Melakukan konfigurasi jaringan sesuai skenario yang diberikan, melakukan pengujian, dan menyusun laporan singkat.}{Hasil konfigurasi jaringan yang berfungsi, laporan ujian, dan screenshot verifikasi.}{Ketepatan konfigurasi jaringan & 10 & Konfigurasi sesuai spesifikasi skenario, semua perangkat terhubung, dan verifikasi konektivitas berhasil penuh. & Konfigurasi sebagian besar benar, sebagian besar perangkat terhubung, tetapi terdapat satu-dua endpoint yang gagal. & Konfigurasi tidak sesuai spesifikasi atau sebagian besar konektivitas gagal. \\ \\
Kualitas subnetting dan penentuan IP & 8 & Subnetting dihitung dengan benar, penugasan IP address konsisten, dan tidak ada konflik address. & Subnetting sebagian besar benar, tetapi terdapat satu kesalahan perhitungan atau penugasan IP yang tidak konsisten. & Subnetting salah atau terdapat banyak konflik IP address. \\ \\
Laporan dan dokumentasi ujian & 2 & Laporan ujian memuat topologi, IP table, langkah konfigurasi, dan bukti verifikasi yang lengkap. & Laporan tersedia tetapi tidak memuat semua komponen atau dokumentasi verifikasi kurang. & Laporan tidak ada atau sangat tidak lengkap. \\}

\assessmentblock{PBL Proyek Praktikum Jaringan Komputer}{Project Based Learning}{SCPMK0504-03301, SCPMK0504-03302, SCPMK0901-03303}{Mahasiswa mengerjakan 12 jobsheet praktikum terintegrasi yang mencakup konfigurasi media jaringan, subnetting, routing, simulasi Packet Tracer, dan analisis traffic jaringan.}{Mengerjakan setiap jobsheet secara sistematis, melakukan konfigurasi dan pengujian, mendokumentasikan seluruh proses, dan mempresentasikan hasil bila diminta.}{12 jobsheet terisi lengkap, laporan per sesi, file simulasi Packet Tracer, dan portofolio praktikum.}{Konsistensi penyelesaian jobsheet & 12 & Semua jobsheet diselesaikan tepat waktu, konfigurasi berhasil, dan dokumentasi lengkap untuk setiap sesi. & Sebagian besar jobsheet diselesaikan, tetapi terdapat beberapa sesi dengan dokumentasi kurang lengkap atau konfigurasi gagal. & Kurang dari setengah jobsheet diselesaikan atau kualitas pengerjaan sangat rendah. \\ \\
Kualitas simulasi Packet Tracer & 8 & Topologi jaringan di Packet Tracer sesuai spesifikasi, semua routing berjalan, dan simulasi traffic jaringan terdokumentasi. & Simulasi Packet Tracer sebagian besar benar, tetapi terdapat routing yang belum berfungsi atau dokumentasi kurang. & Simulasi Packet Tracer tidak sesuai spesifikasi atau tidak dapat mendemonstrasikan konektivitas. \\ \\
Kemampuan analisis traffic dan troubleshooting & 7 & Output ping, traceroute, dan ARP diinterpretasikan dengan benar, masalah diidentifikasi, dan solusi terdokumentasi. & Analisis traffic dilakukan, tetapi interpretasi atau solusi troubleshooting belum lengkap. & Tidak mampu menganalisis output tools jaringan atau mendokumentasikan troubleshooting. \\}

\assessmentblock{UAS Ujian Akhir Praktikum}{UAS Praktikum dan Defense}{SCPMK0504-03301, SCPMK0504-03302, SCPMK0901-03303}{Mahasiswa menyelesaikan ujian akhir praktikum yang menguji kemampuan komprehensif dalam konfigurasi jaringan, routing, dan troubleshooting.}{Mengerjakan skenario konfigurasi jaringan kompleks (termasuk subnetting, routing, dan simulasi), melakukan troubleshooting, dan mempertahankan laporan di hadapan dosen penguji.}{Jaringan yang berfungsi sesuai spesifikasi, laporan akhir, file Packet Tracer, dan presentasi defense.}{Kelengkapan dan ketepatan konfigurasi final & 7 & Seluruh komponen jaringan dikonfigurasi dengan benar, semua skenario routing berjalan, dan konektivitas end-to-end terverifikasi. & Sebagian besar konfigurasi benar, tetapi terdapat satu-dua komponen yang belum berfungsi optimal. & Konfigurasi tidak memenuhi spesifikasi atau banyak komponen gagal diverifikasi. \\ \\
Kemampuan troubleshooting dan analisis & 5 & Masalah pada skenario ujian diidentifikasi dan diselesaikan secara sistematis dengan dokumentasi langkah yang jelas. & Troubleshooting dilakukan tetapi tidak sistematis atau sebagian masalah tidak terselesaikan. & Tidak mampu mengidentifikasi atau menyelesaikan masalah dalam skenario ujian. \\ \\
Defense laporan dan pemahaman konsep & 3 & Mampu menjelaskan setiap keputusan konfigurasi dengan dasar konsep yang kuat dan menjawab pertanyaan dengan tepat. & Dapat menjelaskan keputusan konfigurasi secara umum, tetapi beberapa jawaban belum berdasarkan konsep yang kuat. & Tidak mampu mempertahankan keputusan konfigurasi atau tidak memahami konsep yang diterapkan. \\}

\begin{tblr}{width=\textwidth,colspec={|Q[l,m,3.2cm]|X[l,m]|},hlines,vlines,hspan=minimal,row{1-3}={bg=headgray},rowsep=1.5pt,colsep=3pt}
Jadwal Pelaksanaan: & Minggu 1--17 sesuai RPS. Setiap evaluasi memiliki RTM dan rubrik multi-kriteria. \\
Lain-Lain Yang Diperlukan: & LMS, Cisco Packet Tracer, Wireshark, perangkat jaringan laboratorium, dan laporan praktikum. \\
Pustaka: & Mengacu pustaka utama dan pendukung pada bagian RPS. \\
\end{tblr}
